\chapter{Perspectives et préparation de la réalisation}
\section{Un dossier de conception par sprint}
Après validation du backlog, chaque sprint doit recevoir les mêmes rubriques que les exemples S00 et S01 : fiche de sprint, use case général, fiches détaillées, séquences par cas, classes, activités, maquettes basse fidélité, règles et tests. Les diagrammes ne seront pas copiés à l’identique : seuls les acteurs, objets et interactions utiles au sprint concerné apparaîtront.

\begin{notebox}
\textbf{Principe de traçabilité :} \texttt{US-040} $\rightarrow$ \texttt{S05} $\rightarrow$ \texttt{UC-S05-01} $\rightarrow$ diagramme de séquence $\rightarrow$ classes/écran $\rightarrow$ tests et preuves d’acceptation.
\end{notebox}

\section{Points de cadrage restant ouverts}
\begin{itemize}
\item Sens et périmètre métier de la \textbf{reproduction/saillie} (« tchatel »), et éventuelle rémunération.
\item Politique de \textbf{contact} : messagerie interne, WhatsApp ou combinaison encadrée.
\item Délimitation opérationnelle entre \textbf{adoption gratuite} et \textbf{don}.
\item Processus de \textbf{modération} initiale, pièces justificatives recevables et conservation des informations sensibles.
\item Mise en œuvre du \textbf{multilinguisme}, notamment expressions tunisiennes et orientation RTL.
\item Source et qualité du \textbf{référentiel géographique} tunisien.
\item Conditions de prise en charge des \textbf{espèces dont la détention ou la cession est réglementée}.
\end{itemize}

\section{Trajectoire technique}
\textbf{Étape 1 :} maîtriser le monolithe modulaire, les API REST, les migrations, l’authentification et les tests. \textbf{Étape 2 :} introduire uniquement les outils avancés justifiés par le besoin (Scheduler, Batch, géographie, matching). \textbf{Étape 3 :} étudier l’extraction de microservices après stabilisation des frontières métier et de la chaîne de tests.

\section{Conclusion provisoire}
Le projet dispose d’une vision, d’une étude comparative initiale, d’une solution proposée, d’un backlog structuré et d’une roadmap. Le présent dossier présente des \textbf{modèles de conception préliminaires}. Il sera complété de manière incrémentale à chaque sprint, en distinguant les choix validés des fonctionnalités réellement développées, testées et démontrées.
